% Fragment gotowy do wklejenia do istniejacego skryptu LaTeX.
% W preambule skryptu wymagane:
% \usepackage{tikz}
% \usetikzlibrary{arrows.meta,calc,shadings}
% \usepackage{float} %% dla figure[H]
% Barwy uzytkownika zachowujemy, jesli sa juz zdefiniowane:
\providecolor{manifoldblue}{HTML}{387F9C}
\providecolor{manifoldgold}{HTML}{C1934A}
\providecolor{manifoldred}{HTML}{A85152}
\providecolor{manifoldgreen}{HTML}{438777}

\begin{figure}[H]
\centering
\begin{tikzpicture}[>=Latex,font=\small,line cap=round,line join=round]

 % Warstwy barwy modeluja okragly przekroj rury; wszystkie biegna po tej samej krzywej.
 \def\softtube#1{%
  \draw[manifoldblue!72!black,line width=14.6pt,line cap=butt] #1;
  \foreach \w/\m in {13.2/65,11.8/57,10.4/49,9/41,7.6/34,6.2/28,4.8/23,3.4/19,2/16}{
   \draw[manifoldblue!\m,line width=\w pt,line cap=butt] #1;
  }
  \begin{scope}[xshift=-.35pt,yshift=.7pt]
   \draw[white,opacity=.32,line width=.8pt,line cap=butt] #1;
  \end{scope}
 }
 \pgfdeclareradialshading{chirvolumeball}{\pgfpoint{-8bp}{9bp}}{%
  color(0bp)=(white);color(6bp)=(manifoldblue!8);
  color(13bp)=(manifoldblue!28);color(21bp)=(manifoldblue!52);
  color(25bp)=(manifoldblue!66);color(50bp)=(manifoldblue!70)}
 \def\softball#1#2#3{%
  \shade[shading=chirvolumeball] (#1,#2) circle (#3);
  \draw[manifoldblue!73!black,line width=.75pt] (#1,#2) circle (#3);
 }

 % Nasada jak na szkicu autora: przedni luk ciagly, tylny przerywany.
 \def\softfoot#1#2#3#4#5{%
  \begin{scope}[shift={(#1,#2)},rotate=#3]
   \pgfmathsetmacro{\footradius}{#4/56.9055}
   \begin{scope}
    \clip (-.035,-\footradius)--(0,-\footradius)
     arc[start angle=-90,end angle=90,x radius=#5,y radius=\footradius]
     --(-.035,\footradius)--cycle;
    \softtube{(-.035,0)--(#5,0)}
   \end{scope}
   \draw[manifoldblue!73!black,line width=.65pt]
    (0,-\footradius)
     arc[start angle=-90,end angle=90,x radius=#5,y radius=\footradius];
   \draw[manifoldblue!73!black,opacity=.65,line width=.55pt,
         dash pattern=on 1.2pt off 1.4pt]
    (0,\footradius)
     arc[start angle=90,end angle=270,x radius=#5,y radius=\footradius];
  \end{scope}
 }

 \def\knotpath#1#2{plot[domain=#1:#2,samples=60,smooth,variable=\t]
  ({.52*(sin(\t)+2*sin(2*\t))},
   {.52*(cos(\t)-2*cos(2*\t))})}
 \node[font=\small\bfseries,manifoldblue!75!black] at (0,2.68)
  {(a) Przyklejenie zawęźlonego uchwytu};
 \begin{scope}[xshift=-4.45cm]
  \node[manifoldblue!75!black] at (0,2.10) {$\partial B^3=S^2$};
  \softball{0}{0}{.94}
  \draw[manifoldblue!68,opacity=.30,dashed]
   (-.94,0) arc[start angle=180,end angle=0,x radius=.94,y radius=.27];
  \draw[manifoldblue!68,opacity=.35]
   (-.94,0) arc[start angle=180,end angle=360,x radius=.94,y radius=.27];
  \foreach \x in {-.37,.37}{
   \shade[left color=manifoldgold!20,right color=manifoldgold!62,
          draw=manifoldgold!78!black,line width=.75pt]
    (\x,-.08) ellipse (.155 and .25);
  }
  \node[font=\scriptsize] at (0,-1.40) {dwa dyski $S^0\times D^2$};
 \end{scope}
 \begin{scope}
  \node[manifoldblue!75!black] at (0,2.10) {rdzeń: węzeł $K$};
  \draw[manifoldblue!77!black,line width=2.2pt] \knotpath{0}{360};
  \foreach \a/\b in {8/23,128/143,248/263}{
   \draw[white,line width=6pt] \knotpath{\a}{\b};
   \draw[manifoldblue!77!black,line width=2.2pt] \knotpath{\a}{\b};
  }
  \draw[manifoldgold!80!black,dashed,line width=.8pt]
   (0,-.52) circle (.39);
  \node[manifoldgold!65!black,font=\scriptsize,anchor=north west]
   at (.35,-.66) {$B^3$};
 \end{scope}
 \begin{scope}[xshift=4.60cm]
  \node[manifoldblue!75!black] at (0,2.10) {$T_K=\partial N(K)$};
  \softtube{\knotpath{7}{353}}
  \foreach \a/\b in {8/23,128/143,248/263}{
   \draw[white,line width=19pt,line cap=butt] \knotpath{\a}{\b};
   \softtube{\knotpath{\a}{\b}}
  }
  \shade[shading=chirvolumeball] (0,-.52) circle (.39);
  \draw[manifoldblue!73!black,line width=.75pt]
   ({.39*cos(40)},{-.52+.39*sin(40)})
    arc[start angle=40,end angle=140,radius=.39]
   ({.39*cos(220)},{-.52+.39*sin(220)})
    arc[start angle=220,end angle=320,radius=.39];
  \softtube{\knotpath{7}{12}}
  \softtube{\knotpath{348}{353}}
  \softfoot{-.31498}{-.49298}{-9.85}{14.6}{.10}
  \softfoot{.31498}{-.49298}{189.85}{14.6}{.10}
  \node[font=\scriptsize] at (0,-1.93) {kula i zawęźlony $1$-uchwyt};
 \end{scope}
 \draw[->,manifoldred!80!black,line width=1.1pt] (-3.20,0)--(-1.75,0);
 \draw[->,manifoldred!80!black,line width=1.1pt] (1.75,0)--(2.85,0);
 \node[manifoldred!80!black,font=\scriptsize,align=center]
  at (2.20,.95) {pogrubiamy\\łuk};

 \draw[manifoldblue!20,line width=.45pt] (-6.10,-2.34)--(6.10,-2.34);
 \node[font=\small\bfseries,manifoldblue!75!black] at (0,-2.78)
  {(b) Dodatkowy torus oplatający jedną rurkę};
 \def\greentube#1{%
  \draw[manifoldgreen!65!black,line width=7.2pt,line cap=butt] #1;
  \foreach \w/\m in {6.3/65,5.4/55,4.5/44,3.6/34,2.7/25,1.8/18,.9/12}{
   \draw[manifoldgreen!\m,line width=\w pt,line cap=butt] #1;
  }
 }
 \begin{scope}[xshift=-2.40cm,yshift=-4.85cm,scale=1.30]
  \def\meridian#1#2{plot[domain=#1:#2,samples=48,smooth,variable=\t]
   ({-.65*3/sqrt(34)*sin(\t)},{-.52+.65*cos(\t)})}
  \greentube{\meridian{0}{360}}
  \softtube{\knotpath{0}{360}}
  \foreach \a/\b in {8/23,128/143,248/263}{
   \draw[white,line width=19pt,line cap=butt] \knotpath{\a}{\b};
   \softtube{\knotpath{\a}{\b}}
  }
  \greentube{\meridian{0}{180}}
  \node[manifoldblue!75!black,font=\scriptsize] at (1.82,1.08) {$T_K$};
  \draw[manifoldblue!70!black,line width=.5pt] (1.62,.99)--(1.25,.73);
  \node[manifoldgreen!65!black,font=\scriptsize] at (0,-.90) {$T_\mu$};
 \end{scope}
 \begin{scope}[xshift=3.35cm,yshift=-5.05cm,scale=1.55]
  \node[font=\scriptsize,align=center] at (0,1.12)
   {przekrój płaszczyzną\\normalną do $K$};
  \shade[even odd rule,top color=manifoldgreen!14,
         middle color=manifoldgreen!32,bottom color=manifoldgreen!62]
   (0,0) circle (.75) (0,0) circle (.55);
  \draw[manifoldgreen!65!black,line width=.65pt]
   (0,0) circle (.75) (0,0) circle (.55);
  \fill[manifoldblue!28,draw=manifoldblue!73!black,line width=.65pt]
   (0,0) circle (.20);
  \fill[manifoldblue!73!black] (0,0) circle (.025);
  \node[font=\scriptsize,manifoldblue!75!black] at (0,-1.01)
   {przekrój rurki $N(K)$};
  \draw[manifoldblue!65!black,line width=.5pt] (0,-.86)--(0,-.23);
  \node[font=\scriptsize,manifoldgreen!65!black,anchor=west]
   at (.92,.02) {$N(\mu)$};
  \draw[manifoldgreen!65!black,line width=.5pt] (.89,.02)--(.71,.02);
 \end{scope}

 % (c) Zastepujemy kule torusem; polaczenia jak w panelu (a).
 % Torus ma to samo spektrum blekitow co koniczyna, a przy polaczeniach
 % nie ma gniazd: pozostaja jedynie gladko zaokraglone konce rurek.
 \draw[manifoldblue!20,line width=.45pt] (-6.10,-7.36)--(6.10,-7.36);
 \node[font=\small\bfseries,manifoldblue!75!black] at (0,-7.78)
  {(c) Torus w miejscu kuli przy koniczynie};
 \begin{scope}[shift={(0,-10.10)},scale=1.55]
   \softtube{\knotpath{7}{353}}
   \foreach \a/\b in {8/23,128/143,248/263}{%
      \draw[white,line width=19pt,line cap=butt] \knotpath{\a}{\b};
      \softtube{\knotpath{\a}{\b}}
   }
   % Torus o tym samym typie cieniowania co rurki koniczyny,
   % ale nieco mniejszy, aby poza uchwytami nie stykał się z koniczyną.
   \begin{scope}[shift={(0,-.52)},rotate=-3]
      \draw[manifoldblue!72!black,line width=14.6pt] (0,0) ellipse[x radius=.345,y radius=.365];
      \foreach \w/\m in {13.2/65,11.8/57,10.4/49,9/41,7.6/34,6.2/28,4.8/23,3.4/19,2/16}{%
        \draw[manifoldblue!\m,line width=\w pt] (0,0) ellipse[x radius=.345,y radius=.365];
      }
      \begin{scope}[xshift=-.35pt,yshift=.7pt]
        \draw[white,opacity=.32,line width=.8pt]
          ({.325*cos(142)},{.345*sin(142)}) arc[start angle=142,end angle=45,x radius=.325,y radius=.345];
      \end{scope}
      % Sam otwór bez dodatkowych łuków wewnętrznych.
      \fill[white] (0,0) ellipse[x radius=.118,y radius=.138];
    \end{scope}
   % Polaczenie skopiowane bezposrednio z panelu (a):
   % te same koncowe fragmenty wezla oraz te same nasady softfoot.
   % Nie ma dodatkowych bocznych lacznikow.
   \softtube{\knotpath{7}{12}}
   \softtube{\knotpath{348}{353}}
   % Mniejsze eliptyczne obrzeza: ich zewnetrzna krawedz odpowiada
   % teraz promieniowi rurki (14.6pt / 2), przy skali panelu 1.55.
   \softfoot{-.31498}{-.49298}{-9.85}{9.5}{.10}
   \softfoot{.31498}{-.49298}{189.85}{9.5}{.10}
   \node[anchor=west,text=manifoldblue!72!black,font=\scriptsize]
       at (1.14,-.76) {$T^2$ zamiast $B^3$};
   \draw[manifoldblue!68!black,line width=.55pt]
       (1.10,-.74)--(.36,-.68);
   \node[font=\scriptsize] at (0,-1.51)
       {torus z zawęźlonym $1$-uchwytem};
 \end{scope}

\end{tikzpicture}
\caption{Zawęźlony uchwyt i jego modyfikacje. W (b) zielony torus
meridionalny oplata jedną rurkę węzła, ale pozostaje osobną powierzchnią.
W (c) kula z panelu (a) zostaje zastąpiona torusem w dokładnie
tym samym położeniu: oba końce zawęźlonego $1$-uchwytu dochodzą do nowego korpusu. Po wygładzeniu otrzymujemy
powierzchnię rodzaju dwa.
Różne warstwy w miejscach skrzyżowań oznaczają przejścia nad i pod,
nie zaś przecięcia powierzchni. Schemat topologiczny.}
\label{fig:trefoil-torus-replaces-ball}
\end{figure}
